\section{Introduction}
\label{sec:A-intro}

\subsection{The conjecture}
\label{sec:A-conjecture}

The Quantum Sphaera construction starts from one conjecture: \emph{matter, light and a tachyon
sector are one object, seen from three frames.} In the Companion~\cite{QSC} this is the central
statement \texttt{thm:unity}. The object is a biquaternion kernel built from three complex kernels,
and a projector, the Heisenberg compensator, splits every state into a part that an involution keeps
and a part that it reverses. Matter, light and the tachyon sector are meant to be sectors of this one
kernel, and the frames are meant to be the views in which each sector is seen.

The conjecture has a history that a reader should know. In its first form, ``the same object'' meant
unitary equivalence, and the equivalence was argued from a Stone--von Neumann theorem taken from an
earlier Riemann-hypothesis manuscript of the author. That justification was withdrawn in round
v4.8.0-r04, and the statement has the status of a conjecture since then \Open. What survives today
is more modest and more precise. On the complex three-vector $z=\mathbf E+i\mathbf B$ the Lorentz
group and the electromagnetic duality rotation act, and their orbits are the shells $|s|=r>0$ of the
invariant $s=z\cdot z$, the punctured null cone and the origin (v4.8.0-r07; \rnd{r16}). A matter-type
and a tachyon-type configuration with the same $|s|$ lie on one orbit; light lies on another and is
reached from the first only as a limit (\rnd{r21}). So ``one object'' holds in the sense of an active
symmetry between matter and tachyon \Established, while the physical identification of the sectors
with matter, light and tachyons, and the claim that they are one physical state, remain a \Reading.
Sections~\ref{sec:A-carrier} and~\ref{sec:A-frames} make this precise.

\subsection{What this article does}
\label{sec:A-what}

This article reports the construction as it stands at Companion round v4.10.0-r60. It does not
present a deductive framework. For each result it says which of five kinds it is:
\begin{itemize}
\item \Established: mathematics with a proof, or a computation that anyone can repeat.
\item \Conditional: holds under named choices, settings or assumptions, which are stated with it.
\item \Reading: an interpretation of established mathematics. It may guide work but carries no proof
  weight; \Exploration{} marks results of explorations without a pre-registered criterion.
\item \Open: a question that is posed and not answered.
\item \Refuted: a pre-registered test whose criterion failed. Refuted tests appear only where they ground
  a result of the section in which they stand.
\end{itemize}
The tags are attached in the text where a statement is made. Figure~\ref{fig:A-overview} shows the
construction as a map with the same colours.

\paragraph{How the evidence is organised.}
The Companion grows in \emph{rounds}. A round is appended to the source and never rewritten;
corrections come as new rounds that name what they correct. We cite round v4.10.0-r$NN$ as
\rnd{r$NN$} and round v4.9.0-r$NN$ as \rndold{r$NN$}; rounds of the earlier v4.8.0 line are cited
by their full name, and labels of the Companion main text in typewriter font, e.g.\
(\cite{QSC}, \texttt{prop:compensator}). A test with a decision criterion is written down and committed
as a pre-registered note before the computation, and receives a number H$NN$. When the outcome could
be seen analytically before the computation, the Companion marks the test \emph{not blind}; such a test
checks a computation, not a prediction. Work without a pre-registered criterion is called an
exploration. The bookkeeping of statements, their premises and the guards that block unproved
conclusions is described in a separate article on the QAM language; here we use only its outcome, the
status of each statement.

\paragraph{What is not derived.}
We state this up front because it shapes every later section. The construction derives no particle
mass, no absolute scale and no value of the cosmological constant $\Lambda$. Where a number appears,
such as the frame offset $\bsym\approx0.549$, it is inserted, and we say so. The Koide value $K=2/3$
is not derived: Section~\ref{sec:B-koide} shows that it is equivalent to an overlap $\tfrac12$ between three states,
and that overlap is an input. Against external measurements there is, so far, no confirmed blind
prediction. A review at round \rnd{r05} counted about 20 confirmed, 19 refuted, 9 invalid, 4 open and
28 mixed test outcomes since v4.8.0, 29 of them marked not blind; the confirmed tests without that
mark all check inner mathematics \cite{QSC}. The predictions that are still pending, with their
deadlines, are listed in Section~\ref{sec:open}.

\paragraph{What the construction does supply.}
It supplies exact statements about a carrier and its symmetries: a projector and its image, a sign
rule for three sectors, two field classes under the Lorentz group, duality and scale, the action of
three frame maps, conditions under which three readout states form an octant, and conditional
constructions of one particle generation. It also supplies exclusions: several routes from the
structure to numbers are closed by proofs or refuted tests. We regard both as results.

% F1: overview map of the construction, status at Companion round v4.10.0-r60.
% 2.5D: three stacked layers (carrier, frames and sectors, readouts) drawn as slanted planes with shadows.
% Node fill and edge colour carry the status; a text tag and the line style repeat it for black-and-white print.
\begin{figure}[t]
\centering
\begin{tikzpicture}[
  font=\scriptsize,
  box/.style={draw=black!70, rounded corners=1.2pt, align=center, inner sep=2.2pt, text width=#1},
  box/.default=3.2cm,
  est/.style={fill=st-est!14},
  cond/.style={fill=st-cond!16},
  read/.style={fill=st-read!14},
  open/.style={fill=st-open!16},
  ref/.style={fill=st-ref!14},
  eest/.style={-{Stealth[length=2mm]}, very thick, draw=st-est},
  econd/.style={-{Stealth[length=2mm]}, thick, dashed, draw=st-cond!90!black},
  eopen/.style={-{Stealth[length=2mm]}, thick, densely dotted, draw=st-open},
  eread/.style={-{Stealth[length=2mm]}, thick, dash dot, draw=st-read},
]
% ---------- layer planes (drawn first, back to front) ----------
\newcommand{\qsplane}[3]{% #1 = y offset, #2 = height, #3 = label
  \fill[black!22] (0.15,#1-0.15) -- (13.15,#1-0.15) -- (14.35,#1+#2-0.15) -- (1.35,#1+#2-0.15) -- cycle;
  \filldraw[fill=white, draw=black!55] (0,#1) -- (13,#1) -- (14.2,#1+#2) -- (1.2,#1+#2) -- cycle;
  \node[anchor=south west, font=\footnotesize\bfseries, text=black!65] at (1.25,#1+#2-0.42) {#3};}
\qsplane{7.2}{4.0}{3 \;Readouts}
\qsplane{3.6}{2.6}{2 \;Frames and sectors}
\qsplane{0}{2.6}{1 \;Carrier}

% ---------- layer 1: carrier ----------
\node[box=3.8cm, est] (kern) at (4.0,1.05)
  {biquaternion kernel\\ $q\in\mathbb C\otimes\mathbb H$, $z\in\mathbb C^3$\\ $N(q)=-t^2+s(z)$\quad\textsf{[established]}};
\node[box=3.8cm, est] (comp) at (9.4,1.05)
  {compensator $C=\tfrac12(1+J)$\\ orthogonal projection\\ \textsf{[established]}};

% ---------- layer 2: frames and sectors ----------
\node[box=2.8cm, est] (sect) at (2.9,4.75)
  {sectors by the sign of\\ $s=z_1^2+z_2^2+z_3^2$\\ \textsf{[established]}};
\node[box=2.8cm, est] (maps) at (6.25,4.75)
  {frame maps $z\mapsto iz$,\\ $z\mapsto\bar z$, $J_B$\\ \textsf{[established]}};
\node[box=2.8cm, est] (class) at (9.6,4.75)
  {two field classes:\\ null ($P$) / non-null ($R$)\\ \textsf{[established]}};
\node[box=2.1cm, read] (unity) at (12.75,4.75)
  {one object,\\ three frames\\ \textsf{[conjecture]}};

% ---------- layer 3: readouts ----------
\node[box=3.2cm, cond] (part) at (3.1,8.25)
  {particles from $\mathbb C\otimes\mathbb H\otimes\mathbb O$;\\ Albert algebra, $A_4$\\ \textsf{[conditional]}, Sec.~4};
\node[box=3.2cm, est] (koide) at (6.9,8.25)
  {Koide as geometry: octant\\ $K=\tfrac23\Leftrightarrow$ overlap $\tfrac12$\\ \textsf{[est.]}; $K$ not derived\\ Sec.~5};
\node[box=3.2cm, open] (beta) at (10.7,8.25)
  {$\beta_{\rm sym}\approx0.549$ inserted;\\ counting rule (OP1)\\ \textsf{[open]}, Sec.~6};
\node[box=3.2cm, cond] (act) at (3.9,10.15)
  {common action W1/W2\\ \textsf{[conditional]}, Sec.~7};
\node[box=2.6cm, read] (bang) at (7.6,10.15)
  {Big Bang as transition\\ \textsf{[exploration]}};
\node[box=3.2cm, open] (lam) at (11.4,10.15)
  {$\Lambda$ an input; OP5\\ \textsf{[open]}, Sec.~8};

% ---------- edges ----------
\draw[eest] (kern) -- (sect);
\draw[eest] (kern) -- (maps);
\draw[eest] (comp) -- (class);
\draw[eread] (class) -- (unity);
\draw[eread] (sect.south) to[out=-20,in=-150] (unity.south);
\draw[econd] (kern.west) to[out=160,in=200] (part.west);
\draw[econd] (maps) -- (koide);
\draw[eopen] (maps.north east) -- (beta.south west);
\draw[econd] (act) -- (bang);
\draw[eopen] (act.north) to[out=25,in=155] (lam.north);
\draw[eopen] (beta) -- (lam);
\end{tikzpicture}

\smallskip
\begin{tabular}{@{}l@{\quad}l@{\quad}l@{\quad}l@{}}
\tikz\draw[-{Stealth[length=1.6mm]},very thick,st-est] (0,0)--(0.7,0); established &
\tikz\draw[-{Stealth[length=1.6mm]},thick,dashed,st-cond!90!black] (0,0)--(0.7,0); conditional &
\tikz\draw[-{Stealth[length=1.6mm]},thick,densely dotted,st-open] (0,0)--(0.7,0); open &
\tikz\draw[-{Stealth[length=1.6mm]},thick,dash dot,st-read] (0,0)--(0.7,0); reading
\end{tabular}
\caption{Overview of the construction at Companion round \rnd{r60}. The bottom layer is
the carrier (Section~\ref{sec:A-carrier}), the middle layer the sectors, frame maps and the two field
classes (Section~\ref{sec:A-frames}), the top layer the physical readouts treated in Sections~4--8.
Fill colour, text tag and line style give the status of each node and link. Green: proved
mathematics. Yellow: holds under named choices or settings. Orange: open. Blue: a reading,
conjecture or exploration. The established links into layer~2 are algebra; every arrow into
layer~3 carries a choice or an open step, and no arrow yields a mass, a scale or a value of
$\Lambda$.}
\label{fig:A-overview}
\end{figure}


\subsection{Plan of the article}
\label{sec:A-plan}

Section~\ref{sec:A-carrier} defines the carrier: the biquaternion kernel, the involution $J$ and the
compensator $C=\tfrac12(1+J)$, the sector rule by the sign of $s=z_1^2+z_2^2+z_3^2$, and what $C$
does to a light field. All of it is established mathematics; the names of the sectors are a reading.
Section~\ref{sec:A-frames} introduces the frames M, T, P and R of Convention~F and the maps between
them. The tachyon map $z\mapsto iz$ is the electromagnetic duality rotation, not a Wick rotation and
not time reversal $z\mapsto\bar z$. Under the Lorentz group, duality and scale there are exactly two
classes of non-zero fields, and frame maps, being symmetries, give equalities between frames and no
values. Section~\ref{sec:B-particles} constructs particles from $\mathbb C\otimes\mathbb H\otimes\mathbb O$ under named
choices, up to the family symmetry of the Albert algebra and its breaking to $A_4$. Section~\ref{sec:B-koide} treats
the Koide relation as geometry of three states and shows why it is not a derivation. Section~\ref{sec:beta} is about
the frame offset $\bsym$ and the counting rule behind it, the open problem OP1. Section~\ref{sec:act} presents the
common action W1/W2, which holds under author decisions and supplies no scale, together with a short
exploration of the Big Bang as a transition. Section~\ref{sec:lop5} collects the negative results on $\Lambda$ and
the open problem OP5. Section~\ref{sec:open} lists the open problems and the testable statements with their status,
and Section~\ref{sec:concl} concludes. Appendix~\ref{sec:ledger} gives the status ledger of every statement, Appendix~\ref{sec:repro} the scripts
that reproduce the numbers.
